\documentclass[11pt,a4paper]{article}
\usepackage{fontspec}
\usepackage[english]{babel}
\usepackage{microtype}
\usepackage[margin=2.5cm]{geometry}
\usepackage{amsmath,amssymb}
\usepackage{booktabs,tabularx,array}
\usepackage{enumitem}
\usepackage[hidelinks]{hyperref}
\usepackage{xurl}
\setlist{itemsep=2pt}
\newcommand{\lot}[3]{\href{https://grothendieck.commutator.io/\##1/#2}{shelfmark~#1, p.~#3}}
\newcommand{\q}[1]{“#1”}
\title{Conjectures in Grothendieck's working notes\\[4pt]
\large Dated first appearances, open questions, and what can be formalised}
\author{Michel Hua\\\small\texttt{michel@commutator.io} --- \url{https://grothendieck.commutator.io}}
\date{October 2026}
\begin{document}
\maketitle

\begin{abstract}
\noindent The University of Montpellier holds Alexander Grothendieck's
mathematical working notes, about 16\,000 digitised pages. A transcription
project begun in August 2026 has read more than 11\,000 of them and restated
160 folders in current notation. This note reports what that reading shows
about conjectures. Five well-known conjectures appear in the notes, some
earlier than the texts usually cited for them: the standard conjectures
(1965), the de Rham--Betti form of the period conjecture (1965), the homotopy
hypothesis (1975), the section conjecture, and the programme behind the
Grothendieck--Teichmüller group. The notes also pose questions that have no
printed source. For each item we give the place and date in the fonds, the
published source to cite, its status, whether it is already in the
\emph{Formal Conjectures} repository, and whether current Mathlib can state
it. Everything rests on transcriptions that no person has yet read in full
against the manuscripts; Section~\ref{sec:caveats} says what follows from that.
\end{abstract}

\section{The corpus and how it was read}

The fonds is the deposit of Grothendieck's mathematical notes at Montpellier,
catalogued by folder (« cote ») and published as facsimiles at
\url{https://grothendieck.umontpellier.fr}. The project at
\url{https://grothendieck.commutator.io} works on it in four layers.

\begin{itemize}
\item \textbf{Transcription.} Each batch of twenty pages is transcribed into
  LaTeX with a critical apparatus (illegible, doubtful, struck, added), by
  Anthropic's Claude models (Opus~5 and Opus~5.5) under written editorial
  rules. More than 11\,000 pages are done.
\item \textbf{Modernised readings.} 160 folders are restated whole, in French,
  in current notation, with footnotes marking every departure from the page.
  A reading is an interpretation; where it disagrees with the transcription,
  the transcription is the record.
\item \textbf{Findings.} A register of 458 entries lists statements of the
  folders that may not be in the literature, each written as a claim about the
  literature, with what was searched and what would settle it. Its status
  vocabulary is used below: \emph{matched} (found in print), \emph{candidate}
  (searched, not found), \emph{unsearched}, \emph{refuted}.
\item \textbf{Lean.} An appendix proves in Lean~4, on Mathlib, 1\,018
  theorems drawn from the readings, using only Lean's standard axioms
  (\url{https://grothendieck.commutator.io/article/grothendieck-lean.pdf}).
  Lean checks the proofs, not that a statement is faithful to the page.
\end{itemize}

References of the form \q{shelfmark~$N$, p.~$k$} below point to the reading
view of the folder, which shows the transcription beside Montpellier's
facsimile.

\section{Five conjectures and where they first appear}

Table~\ref{tab:five} sums up this section. The dates are those of the leaves or
of Montpellier's inventory, whose brackets mark a date the archivists
inferred.

\begin{table}[h]
\small
\begin{tabularx}{\textwidth}{>{\raggedright}p{2.7cm} >{\raggedright}p{3.2cm} >{\raggedright}X >{\raggedright\arraybackslash}p{2.6cm}}
\toprule
Conjecture & In the fonds & Published source to cite & Mathlib today \\
\midrule
Standard conjectures & 14, p.~44: letter to Serre, 27 Aug.\ 1965; 16 (1966--67) & Grothendieck, Bombay 1968 (1969); Kleiman 1968, 1994 & no Weil cohomology \\
Period conjecture (de Rham--Betti) & 13, pp.~47--48 (1965) & Grothendieck 1966; André 2004; Bost--Charles 2016 & elliptic case statable \\
Homotopy hypothesis & 134-2, p.~42: letter to Breen, 5 Feb.\ 1975 & \emph{Pursuing Stacks} (1983); Maltsiniotis 2010 & partly \\
Section conjecture & 149, p.~65 (undated) & letter to Faltings, 27 June 1983 (1997) & no étale $\pi_1$ \\
Galois action and $\widehat{GT}$ & 119 (\emph{Esquisse}, 1984); 142 ([from 1978]); 147 & \emph{Esquisse} (1997); Drinfeld 1990; Ihara 1990 & partly \\
\bottomrule
\end{tabularx}
\caption{The five conjectures. None of them is in \emph{Formal Conjectures} as of 17 October 2026 (code search on the repository).}
\label{tab:five}
\end{table}

\subsection{The standard conjectures (1965)}
On 27 August 1965 a letter to Serre states the standard conjectures, among
them the positivity conjecture of Hodge type (\lot{14}{3}{44}); the letter is
printed in the \emph{Correspondance Grothendieck--Serre}. Folder~16
(1966--67) works out their algebraic side: Lefschetz operators, the
pseudo-inverse $\Lambda$, Poincaré--Hodge algebras. Part of that algebra is
proved in the Lean appendix (the normalisation of $\mathfrak{sl}_2$ and the
pseudo-inverse, folder~16), and the register records one statement of the
folder as possibly unpublished: the Lefschetz-type conjecture follows from
$D(X)$ and the existence of \emph{some} algebraic isomorphism
$H^{2n-i} \to H^i$, not necessarily inverse to $L^{n-i}$ (unsearched).

\emph{Status.} Open in general: the Hodge-type conjecture holds in
characteristic~0 by the Hodge--Riemann relations; the Lefschetz-type
conjecture is known in dimension at most~2 and for abelian varieties.

\subsection{The period conjecture (1965)}
Folder~13, \q{Groupe de Galois motivique} (1965), asks on pp.~47--48 whether,
over an algebraic extension $k$ of $\mathbb{Q}$, the functor sending a motive
to its de Rham realisation, its Betti realisation and the comparison
isomorphism between them is fully faithful. It reduces the question to
$k = \overline{\mathbb Q}$ and recasts it as $V^G = V \cap g(\overline V)$ for
the motivic Galois group, noting that the rank-one (Tate) case works. This is
the de Rham--Betti form of what is now called Grothendieck's period conjecture
(register: \emph{matched}, Bost--Charles 2016, Conj.~2.12).

\emph{Status.} Open. The elliptic case can be stated with Mathlib's
\texttt{PeriodPair} and \texttt{weierstrassP}: for a lattice with algebraic
invariants and no complex multiplication, the periods and quasi-periods
$\omega_1, \omega_2, \eta_1, \eta_2$ should have transcendence degree~4; with
complex multiplication the degree is~2 (Chudnovsky, 1976).

\subsection{The homotopy hypothesis (1975)}
On 5 February 1975, the postscript of a letter to Lawrence Breen sets out the
\q{yoga} by which an $n$-category in groupoids, up to $n$-equivalence,
\q{est essentiellement la même chose} as a semi-simplicial set up to homotopy
whose $\pi_i$ are neglected for $i \geq n+1$ (\lot{134-2}{3}{42}). It is the
homotopy hypothesis in its first form; \emph{Pursuing Stacks} takes it up in
1983 for non-strict $\infty$-groupoids. The fonds also holds the handwritten
French original of the third letter to Breen (17--19 July 1975,
shelfmark~134-2, pp.~145--184), of which the folder otherwise holds only English versions, one of them
R.~Brown's translation.

\emph{Status.} Depends on the model: false for strict $n$-groupoids from
$n = 3$ (Simpson), regarded as open for Grothendieck's own globular
$\infty$-groupoids in Maltsiniotis's formulation. An OpenAI preprint dated 24 September 2026, not peer-reviewed, announces a proof for every Grothendieck coherator in the Ara--Henry convention, together with Henry's pushout conjecture on elementary expansions (OpenAI, \q{The Grothendieck homotopy hypothesis via elementary expansions}).

\subsection{The section conjecture}
An undated page of folder~149 concludes \q{Nos conjectures impliquent donc
qu'il y a exactement deux sortes de telles sections}, those of rational points
and those of points at infinity (\lot{149}{4}{65}). This is the section
conjecture in its affine form, usually cited from the letter to Faltings of
27 June 1983.

\emph{Status.} Open; injectivity is known, and so are a birational $p$-adic
variant (Koenigsmann) and the real analogue.

\subsection{The Galois action on Teichmüller towers}
The \emph{Esquisse d'un programme} (1984; folder~119) is at the origin of the
group Drinfeld named after Grothendieck and Teichmüller. Folders~142,
\q{Action de Galois sur Teichmüller} ([from 1978]), and 147, \q{Structure à
l'infini des $M_{g,\nu}$}, work on the same subject. Folder~142 lists
necessary conditions for an automorphism of the profinite fundamental groupoid
of a curve to come from $\mathrm{Gal}(\overline K/K)$ (pp.~94--101; register:
\emph{candidate}).

\emph{Status.} Ihara's question --- is $\mathrm{Gal}(\overline{\mathbb Q}/\mathbb Q)
\to \widehat{GT}$ an isomorphism? --- is open; the map is injective by Belyi's
theorem.

\section{Questions posed in the notes, with no printed source}

The register also holds questions that the notes raise and that the project
has not found in print. They cannot go into \emph{Formal Conjectures} as they
stand, since that repository asks for a published source; they are listed so
that a reader who knows the literature can say whether they are known.

\begin{itemize}
\item \textbf{The 27 lines around a triangle} (folder~69, \q{Graphes
  cubiques}, [from 1976], pp.~7--8, 14--15, 25--26, 113; unsearched). The
  intersection graph of the 27 lines with a marked triangle $t$ is described as
  the data of a 3-element set, an $\mathbb F_2$-plane $V$, and a torsor under
  the Baer product of three pull-backs of the extension of $V$ by
  $\mathbb F_2$ that splits over no subgroup of order~2 --- the quaternion
  group $Q_8$; $\operatorname{Aut}(C, t)$ has order $1152 = |W(F_4)|$. The
  group is classical; the torsor reconstruction may not be. A general version
  for graphs whose edges lie in unique triangles (folder~74) is proved in Lean.
\item \textbf{Cohomological dimension and imperfection} (folder~31, SGA~4
  residues, [1963--1973], p.~84; unsearched). The bound on the
  $\ell$-cohomological dimension of affine schemes in characteristic $\ell$
  would be governed by the degree of imperfection $\nu(k)$,
  $\ell^{\nu(k)} = [k : k^\ell]$, rather than by the dimension.
\end{itemize}

Some questions of the notes are now settled, and serve as a check on the
method: whether the monodromy of projective structures on a fixed Riemann
surface is injective (folder~159; yes, Poincaré); whether sequential colimits
of presheaves on a test category are homotopy colimits (folder~107; not in
general, Milnor's $\lim^1$); whether stratifying cohomology is invariant under
nilpotent immersions (folder~7; no, as the printed exposé says); and in
folder~68, \q{Jeux de position} ([from 1978]), a statement left under a
question mark --- that winning positions are reached by finitely many
backward steps --- which needs transfinite iteration, $\omega + 1$ already
for a position with infinitely many successors.

\section{Relation to recent machine-written proofs}
On 6 October 2026 OpenAI made public a repository of 719 manuscripts
produced by an internal model. Four families announce proofs of statements
from Grothendieck's programme or bearing his name: the homotopy hypothesis
for weak globular $\infty$-groupoids; the section conjecture, local
$p$-adic and over $\mathbb Q$ for the modular curves $X_0(N)$ and $X_1(N)$ of
genus at least two; the Deligne--Drinfeld conjecture on the
Grothendieck--Teichmüller Lie algebra; and the rational Hodge conjecture for
abelian varieties with complex multiplication, hence the Hodge-type standard
conjecture for them. None has been peer-reviewed, and only two are partly
formalised in Lean. This note does not assess them. The fonds supplies dated
points to which a confirmed proof can be related, and it does not attribute
those conjectures: Deligne--Drinfeld is Deligne's and Drinfeld's, the Hodge
conjecture is not Grothendieck's, and the proved forms of the homotopy
hypothesis come from later work.

\section{Formalisation}

\paragraph{What exists.} The Lean appendix proves 1\,018 theorems from the
readings, on a Mathlib master of October 2026; general lemmas from it are
proposed to Mathlib as pull requests. One overlap with \emph{Formal
Conjectures} was found: a proof in folder~36 (SGA~7, [about 1967--1973], on invariant cycles in
$H^1$ over a henselian trait) holds \q{at least modulo resolution of
singularities}, whose solved and open cases are stated in
\texttt{Wikipedia/ResolutionOfSingularities.lean}; the appendix now cites it.

\paragraph{What we propose.} In order of feasibility:
\begin{enumerate}
\item the period conjecture, elliptic case, as \texttt{research open}, with
  the CM case as \texttt{research solved}; it needs only the Weierstrass zeta
  function and the quasi-periods on top of Mathlib's \texttt{PeriodPair};
\item a weak form of Ihara's question for the group $\widehat{GT}_0$ defined
  by Drinfeld's relations (I) and (II) on the profinite free group $\widehat F_2$;
\item the section conjecture and the standard conjectures as issues only,
  with precise informal statements, until étale fundamental groups and Weil
  cohomology exist in Mathlib.
\end{enumerate}
Each would cite the published source of Table~\ref{tab:five}, with the fonds
as context.

\section{Caveats and how to cite}\label{sec:caveats}
\begin{itemize}
\item \textbf{Unchecked readings.} The transcriptions are first passes by a
  language model. A few pages have been read against the facsimile by a
  person; none of the folders above has been checked whole. Quotations are
  given as transcribed; the facsimile is the authority.
\item \textbf{No priority claims.} A conjecture appearing in a folder says
  where Grothendieck wrote it, not that he was first; dates in brackets are
  the archivists' inferences.
\item \textbf{Rights.} The fonds is under copyright; publishing transcriptions
  requires the heirs' authorisation, which is being sought. Statements of
  mathematics are cited here; the manuscripts are not reproduced.
\item \textbf{AI use.} The transcriptions, readings, register, Lean proofs and
  this note were produced with Anthropic's Claude models under the author's
  direction; the author checked and takes responsibility for the content.
\end{itemize}

\section*{References}
\small\raggedright
\begin{itemize}[leftmargin=*]
\item Y. André, \emph{Une introduction aux motifs}, Panoramas et Synthèses 17, SMF, 2004.
\item J.-B. Bost and F. Charles, \q{Some remarks concerning the Grothendieck period conjecture}, \emph{J. reine angew. Math.} 714 (2016).
\item P. Colmez and J.-P. Serre (eds.), \emph{Correspondance Grothendieck--Serre}, SMF, 2001.
\item V. G. Drinfeld, \q{On quasitriangular quasi-Hopf algebras and on a group that is closely connected with $\mathrm{Gal}(\overline{\mathbb Q}/\mathbb Q)$}, \emph{Leningrad Math. J.} 2 (1991).
\item Google DeepMind, \emph{Formal Conjectures}, \url{https://github.com/google-deepmind/formal-conjectures}.
\item A. Grothendieck, \q{On the de Rham cohomology of algebraic varieties}, \emph{Publ. Math. IHÉS} 29 (1966).
\item A. Grothendieck, \q{Standard conjectures on algebraic cycles}, in \emph{Algebraic Geometry (Bombay Colloquium, 1968)}, OUP, 1969.
\item A. Grothendieck, \q{Esquisse d'un programme} and letter to G. Faltings (27 June 1983), in L. Schneps and P. Lochak (eds.), \emph{Geometric Galois Actions} 1, LMS Lecture Notes 242, CUP, 1997.
\item A. Grothendieck, \emph{À la poursuite des champs}, vol.~I, ed. G. Maltsiniotis, Documents mathématiques 20, SMF, 2022.
\item Y. Ihara, \q{Braids, Galois groups, and some arithmetic functions}, \emph{Proc. ICM Kyoto} (1990).
\item S. Kleiman, \q{Algebraic cycles and the Weil conjectures}, in \emph{Dix exposés sur la cohomologie des schémas}, 1968; \q{The standard conjectures}, \emph{Proc. Symp. Pure Math.} 55 (1994).
\item G. Maltsiniotis, \q{Grothendieck $\infty$-groupoids, and still another definition of $\infty$-categories}, arXiv:1009.2331 (2010).
\item OpenAI, \q{The Grothendieck homotopy hypothesis via elementary expansions}, preprint, 24 September 2026, \url{https://github.com/openai/math/blob/main/preprints/The-Grothendieck-homotopy-hypothesis-via-elementary-expansions-September-24-2026/paper.pdf}.
\item OpenAI, \emph{math}: mathematical manuscripts and supporting proof artifacts produced by an internal OpenAI model, \url{https://github.com/openai/math}, 6 October 2026.
\item J. Stix, \emph{Rational Points and Arithmetic of Fundamental Groups}, LNM 2054, Springer, 2013.
\item Université de Montpellier, fonds Alexandre Grothendieck, \url{https://grothendieck.umontpellier.fr}.
\end{itemize}

\end{document}
